\chapter{Análisis y discusión}

Los ensayos demostraron que una trayectoria puede producir avance repetible y
mantener inclinaciones acotadas sin ejecutar el patrón de apoyo supuesto. Esta
diferencia es relevante porque el desplazamiento visual, por sí solo, habría
permitido aceptar incorrectamente el gateo. La publicación sincronizada de
fase y contacto permitió contrastar la intención del controlador con la
interacción pie--suelo.

El debounce no mejoró mecánicamente la marcha. Su función fue exigir
persistencia antes de declarar vuelo. Por ello redujo la coincidencia temporal
global y añadió aproximadamente el retardo configurado a los despegues
delanteros. En las patas traseras, los eventos crudos duraron alrededor de
0,074~s y desaparecieron al aplicar el criterio de 0,12~s. La interpretación
defendible es que estos eventos correspondieron a interrupciones breves del
contacto de Gazebo, no a una fase de vuelo sostenida.

La corrección del temporizador sí produjo una mejora verificable de ejecución:
la cadencia pasó de ciclos cercanos a 4,80~s a ciclos de 4,32~s coherentes con
la configuración. Debido a que esta modificación cambió la velocidad real, las
bolsas anteriores y posteriores se analizaron como versiones experimentales
separadas.

La comparación Gazebo--MuJoCo también tuvo un alcance diferente. Gazebo aportó
pose corporal, contacto y estabilidad, mientras que el adaptador de MuJoCo
empleado permitió evaluar cadencia y articulaciones. No se infirió equivalencia
de avance o estabilidad entre simuladores a partir de variables no observadas.

Los márgenes de estabilidad y las posiciones de pie dependen todavía de una
geometría y centro de masa nominales. Asimismo, los modelos de contacto no
representan directamente los futuros sensores físicos. Los resultados de
simulación respaldan el desarrollo del controlador y el protocolo, pero no
sustituyen la validación del robot real.

La capa correctiva se entrenó con cinco semillas y se conectó a Gazebo, pero
ninguna escala residual positiva cumplió simultáneamente los criterios de
avance y postura. La escala cero reprodujo el nominal, lo que descartó un fallo
del bypass, mientras las escalas positivas mostraron interferencia con la
marcha. En consecuencia, el resultado válido es negativo: la política se
rechaza para transferencia. La comparación fue descriptiva y no completamente
emparejada; no sustituye la campaña final prevista después de rediseñar el
entrenamiento.

Las pruebas provocadas demostraron que el supervisor recibe estímulos de pose,
contacto, estabilidad y trayectoria y ordena \texttt{stand}. Sin embargo,
\texttt{stand} es una orden lógica, no un corte de energía. En hardware deberá
existir una ruta independiente por OE, además de un rearme explícito y una
medición de falsos positivos antes de habilitar por defecto los watchdogs de
contacto, margen y telemetría.
